% !TeX root = ../../Book3.tex
% 纲要标题：### 17.2 投射维数
% 纲要位置：工作记录/计划纲要.md，第 2526 行。
% 纲要细目（写作范围）：
% * 有限生成模的投射维数
% * Ext、Tor 与投射维数判据
% * 维数平移及一般投射模的 Ext 判据在此证明；第四卷第9.5节与第11章作系统整理
% * 剩余域上界、正则元换环及极小消解的同调对照表
% * 测试模选择的 Tor、Ext 反例及完备化保持 Betti 数、投射维数

\section{投射维数}
\label{b3:sec:1702}

一个模有无限自由消解，并不表示每种消解都必须无限长。投射维数问最短长度，而 Tor 消失与极小消解的末项提供可以实际检测这一长度的条件。

% 纲要内容（仅作写作计划，尚非教材正文）：
%
% * 有限生成模的投射维数
% * Ext、Tor 与投射维数判据
% * 维数平移及一般投射模的 Ext 判据在此证明；第四卷第9.5节与第11章作系统整理
% * 剩余域上界、正则元换环及极小消解的同调对照表
%

\subsection{有限生成模的投射维数}
\label{b3:subsec:1702:01}

\begin{definition}\label{b3:def:projective-dimension}
设 \(R\) 为交换环，\(M\ne0\) 为 \(R\) 模。若存在正合列
\[
0\longrightarrow P_n\longrightarrow\cdots\longrightarrow P_0
\longrightarrow M\longrightarrow0
\]
且各 \(P_i\) 投射，则使这种列存在的最小非负整数 \(n\) 称为 \(M\) 的\term[toushe-weishu]{投射维数}[projective dimension]，记 \(\operatorname{pd}_RM\)。若不存在，记为 \(\infty\)；约定 \(\operatorname{pd}_R0=-\infty\)。
\end{definition}

定义中的 \(P_i\) 只要求投射，并不预设自由或有限生成。对 Noether 局部环上的有限生成模，极小自由消解的各项虽都有限自由，却已经实现这类投射消解的最短长度。投射维数衡量消解必须延伸多少步；任意消解却可以插入 \(R\xrightarrow{1}R\)，人为增加长度。

\begin{proposition}\label{b3:prop:pd-minimal-length}
设 \((R,\mathfrak m,k)\) 为交换 Noether 局部环，\(M\ne0\) 为有限生成的 \(R\) 模。若 \(F_\bullet\rightarrow M\) 是极小自由消解，则
\[
\operatorname{pd}_RM=\sup\{\,i:F_i\ne0\,\}
=\sup\{\,i:\operatorname{Tor}_i^R(M,k)\ne0\,\}.
\]
\end{proposition}
\begin{proof}
投射模的提升性质使自由消解的比较映射论证同样适用于投射消解。因此有限投射消解在高于其长度的次数计算出零 Tor，故极小消解的这些项由\cref{b3:prop:minimal-resolution-betti}为零。反向，若极小消解在次数 \(r\) 的项非零，而 \(F_{r+1}=0\)，则全部后续项也为零，它本身就是长度 \(r\) 的投射消解；前一方向又排除了任何更短的投射消解，故两个上确界都为 \(r\)。若每个 \(F_i\) 都非零，则 Tor 在任意高次数均非零，不可能另有有限投射消解，两个上确界及投射维数都为 \(\infty\)。
\end{proof}

\subsection{Ext、Tor 与投射维数判据}
\label{b3:subsec:1702:02}

极小消解的微分在剩余域 \(k\) 上全为零，所以 \(\operatorname{Tor}_i^R(M,k)\) 和 \(\operatorname{Ext}_R^i(M,k)\) 直接记录 \(F_i\) 的秩。任一项同调为零就迫使该自由项为零，\cref{b3:cor:betti-gap-termination}再使所有后续项为零。因此在这里，仅测试剩余域便能判断消解是否终止，无须逐个测试所有模。

\begin{theorem}\label{b3:thm:local-pd-vanishing}
设 \((R,\mathfrak m,k)\) 为交换 Noether 局部环，\(M\) 为有限生成的 \(R\) 模，\(n\ge0\) 为整数。下列条件等价：
\begin{equivalences}
\item \(\operatorname{pd}_RM\le n\)；
\item \(\operatorname{Tor}_{n+1}^R(M,k)=0\)；
\item \(\operatorname{Ext}_R^{n+1}(M,k)=0\)；
\item 对每个 \(R\) 模 \(N\)，有 \(\operatorname{Ext}_R^{n+1}(M,N)=0\)。
\end{equivalences}
若 \(M\ne0\) 且 \(r=\operatorname{pd}_RM<\infty\)，则每个非零有限生成的 \(R\) 模 \(N\) 都满足 \(\operatorname{Ext}_R^r(M,N)\ne0\)。
\end{theorem}
\begin{proof}
\(M=0\) 时所有条件直接成立。设 \(M\ne0\)，前面三项分别等价于极小消解的 \(F_{n+1}=0\)，由\cref{b3:cor:betti-gap-termination}得到同一终止条件。第一项使 Hom 复形在次数 \(n+1\) 后为零，推出第四项；第四项取 \(N=k\) 推出第三项。

若最后非零项为 \(F_r\)，约定 \(r=0\) 时下面的像为零，便可统一写成
\[
\operatorname{Ext}_R^r(M,N)
=\operatorname{Hom}_R(F_r,N)/\operatorname{im}(\mathrm{d}_r^*).
\]
极小性给 \(\operatorname{im}(\mathrm{d}_r^*)\subseteq\mathfrak m\operatorname{Hom}_R(F_r,N)\)。记 \(T=\operatorname{Hom}_R(F_r,N)\)，它是非零有限个 \(N\) 的直和，所以非零且有限生成。对局部环 \(R\) 上的 \(T\) 用 Nakayama 引理，得到 \(T/\mathfrak mT\ne0\)。所列 Ext 商仍满射到 \(T/\mathfrak mT\)，故非零；这里必须保留 \(N\) 非零且有限生成的条件。\(r=0\) 时像为零，同样成立。
\end{proof}

任意其他测试模未必能看见这些自由项；\cref{b3:prop:pd-test-module-counterexamples}将给出 Tor 的测试模换成一个非零商、Ext 的测试模换成环本身后判据失效的计算。

\ProseParagraphBreak
对非零有限生成局部模，深度取 \(\operatorname{Ext}_R^i(k,M)\) 的首次非零次数，投射维数则取 \(\operatorname{Ext}_R^i(M,k)\) 的非零次数上确界；后者可能为无穷。两个变量和两个端点都不同。在 \(R=k[X,Y]_{(X,Y)}\) 上，即使两边都取同一个模 \(M=k\)，所得数值仍不同。用 \(X,Y\) 的 Koszul 消解映入 \(k\)，微分全为零，零、一、二次 Ext 的 \(k\) 维数分别为 \(1,2,1\)，其余为零。因此
\[
\begin{array}{c|c|c}
\text{检测目标}&\text{Ext 计算}&\text{所得数值}\\\hline
\operatorname{depth}_Rk&\operatorname{Ext}_R^0(k,k)=k,
\ \text{首次非零次数为 }0&0\\
\operatorname{pd}_Rk&\operatorname{Ext}_R^2(k,k)=k,
\ \operatorname{Ext}_R^{i>2}(k,k)=0&2
\end{array}
\]
同一组 Ext 的首次非零和最后非零分别给出深度与投射维数；把两个端点混为一谈，也会在这个例子中产生错误。

\subsection{维数平移}
\label{b3:subsec:1702:03}

这里的“维数平移”指同调次数的平移：从一个模换到它的关系模，Ext 和 Tor 的次数相应改变一，与 Krull 维数 \(\dim R\) 无关。

\begin{proposition}\label{b3:prop:dimension-shifting-local}
设 \(R\) 为交换环，\(0\rightarrow K\rightarrow P\rightarrow M\rightarrow0\) 正合且 \(P\) 投射。对任意 \(R\) 模 \(N\) 和 \(i\ge1\)，有自然同构
\[
\operatorname{Ext}_R^i(K,N)\cong\operatorname{Ext}_R^{i+1}(M,N),
\qquad
\operatorname{Tor}_i^R(K,N)\cong\operatorname{Tor}_{i+1}^R(M,N).
\]
\end{proposition}
\begin{proof}
投射模是自由模的直和项，故它的所有正次 Ext 和 Tor 为零。在\cref{b3:lem:ext-basic-properties}和\cref{b3:lem:tor-basic-properties}的长正合列中，连接映射两边相邻的 \(P\) 项都为零，因而连接映射为同构。
\end{proof}

\begin{proposition}[投射模的 Ext 判据]\label{b3:prop:projective-ext-criterion}
设 \(R\) 为交换环，\(P\) 为 \(R\) 模。以下条件等价：
\begin{equivalences}
\item \(P\) 投射；
\item 对每个 \(R\) 模 \(N\)，有 \(\operatorname{Ext}_R^1(P,N)=0\)。
\end{equivalences}
\end{proposition}
\begin{proof}
若 \(P\) 投射，则 \(P\) 是某个自由模 \(F\) 的直和项。Ext 在第一变量保持有限直和，而自由模的正次 Ext 为零，因此 \(\operatorname{Ext}_R^1(P,N)=0\) 对每个 \(N\) 成立。

反过来，取自由模满射 \(q:F\rightarrow P\)，令 \(K=\ker q\)。要把 \(P\) 实现为自由模的直和项，只需找到截面 \(s:P\to F\)，使 \(qs=\operatorname{id}_P\)；这就是把恒等映射沿 \(q\) 提升。对短正合列 \(0\rightarrow K\rightarrow F\xrightarrow{q}P\rightarrow0\) 应用第二变量的 Ext 长正合列，得到
\[
\operatorname{Hom}_R(P,F)\xrightarrow{q\circ-}\operatorname{Hom}_R(P,P)
\longrightarrow\operatorname{Ext}_R^1(P,K)=0.
\]
因此所需的 \(s\) 存在。每个 \(u\in F\) 都可写成 \(sq(u)+(u-sq(u))\)，第二项属于 \(K\)，而 \(s(P)\cap K=0\)，所以 \(F=s(P)\oplus K\)。原短正合列分裂，\(P\) 是自由模 \(F\) 的直和项，故投射。
\end{proof}

设 \((R,\mathfrak m)\) 为交换 Noether 局部环，\(M\) 为有限生成的 \(R\) 模。令 \(\Omega^0M=M\)，并令 \(\Omega^{j+1}M\) 为极小满射 \(F_j\rightarrow\Omega^jM\) 的核；这些是\cref{b3:def:syzygy-local}中合冲模在极小消解上的选择。沿消解换到下一层合冲模，每次便将相应 Ext 的次数降低一。迭代 \(n\) 次后，对 \(n\ge0\) 与任意 \(R\) 模 \(N\)，有
\[
\operatorname{Ext}_R^1(\Omega^nM,N)\cong\operatorname{Ext}_R^{n+1}(M,N).
\]

特别地，若 \(M\) 的投射维数为有限正整数 \(r\)，则 \(\Omega M\ne0\) 且
\(\operatorname{pd}_R\Omega M=r-1\)。极小消解去掉零次项就是 \(\Omega M\) 的极小消解，所以这个等式还可以不经一般判据而直接看出。

\subsection{剩余域上界与正则元换环}
\label{b3:subsec:1702:04}

\begin{corollary}\label{b3:cor:residue-pd-bound}
设 \((R,\mathfrak m,k)\) 为交换 Noether 局部环。每个有限生成的 \(R\) 模 \(M\) 都满足
\(\operatorname{pd}_RM\le\operatorname{pd}_Rk\)。有限生成模的直和因子具有不大于原模的投射维数。
\end{corollary}
\begin{proof}
若 \(k\) 有长度 \(n\) 的投射消解，用它计算第二变量便得
\(\operatorname{Tor}_{n+1}^R(M,k)=0\)，由\cref{b3:thm:local-pd-vanishing}得到结论。若 \(\operatorname{pd}_Rk=\infty\)，不等式没有额外内容。对直和因子，张量复形及其同调保留有限直和；原模的该阶 Tor 为零时，各因子也为零，再应用同一判据。
\end{proof}

因此，只要剩余域的投射维数有限，它就给出本环所有有限生成模的投射维数的统一上界。

\ProseParagraphBreak
对消解逐项模掉 \(x\)，就是张量 \(R/(x)\)，而张量一般不能保留左端正合性。下面要求 \(x\) 同时在环和系数模上为非零因子，正是为了保证取商后的增广消解仍正合。

\begin{lemma}\label{b3:lem:regular-element-resolution-basechange}
设 \((R,\mathfrak m)\) 为交换 Noether 局部环，\(M\ne0\) 为有限生成的 \(R\) 模，\(x\in\mathfrak m\) 在 \(R\) 和 \(M\) 上均为非零因子。令 \(B=R/(x)\)。把 \(M\) 的极小自由消解模掉 \(x\)，得到 \(M/xM\) 的极小 \(B\) 自由消解。因此
\[
\operatorname{pd}_B(M/xM)=\operatorname{pd}_RM.
\]
\end{lemma}
\begin{proof}
因为 \(x\) 在 \(R\) 上为非零因子，\(B\) 有长度一的 \(R\) 自由消解
\[
0\longrightarrow R\xrightarrow{x}R\longrightarrow B\longrightarrow0.
\]
用这条消解及\cref{b3:lem:tor-balanced-computation}计算，得到
\[
\operatorname{Tor}_1^R(M,B)\cong\ker(x:M\rightarrow M)=0,
\qquad
\operatorname{Tor}_i^R(M,B)=0\quad(i\ge2).
\]
一次同调的消失用到了 \(x\) 在 \(M\) 上为非零因子；高次消失则来自这条消解没有更高项。因此，若 \(F\) 为 \(M\) 的极小自由消解，便有
\[
H_i(F\otimes_RB)=\operatorname{Tor}_i^R(M,B)=0\quad(i>0),
\qquad H_0(F\otimes_RB)=M\otimes_RB\cong M/xM.
\]
所以模掉 \(x\) 后的增广复形仍正合。各项是有限秩自由 \(B\) 模，微分矩阵的元素仍落在 \(B\) 的极大理想 \(\mathfrak m/(x)\) 中，故消解仍极小。Nakayama 引理给出 \(M/xM\ne0\)，而 \(F_i/xF_i\ne0\) 当且仅当 \(F_i\ne0\)。由\cref{b3:prop:pd-minimal-length}，两个极小消解的长度相同，包括无穷情形，得到所述等式。
\end{proof}

仍设 \((R,\mathfrak m)\) 为交换 Noether 局部环，\(x\in\mathfrak m\) 在 \(R\) 上为非零因子，若改取 \(M=B=R/(x)\)，则 \(xM=0\)。长度一的 \(R\) 自由消解的微分为 \(x\)，故它极小，给出 \(\operatorname{pd}_RB=1\)；而 \(B\) 在自身上自由，\(\operatorname{pd}_BB=0\)。此时 \(\operatorname{Tor}_1^R(B,B)=B\ne0\)，所以模掉 \(x\) 后的消解不再正合。共同非零因子的条件在这里不能删去。

\ProseParagraphBreak
自由项的秩、剩余域上的同调与消解的终止性，现在都可以在极小消解中读取。\cref{b3:prop:minimal-resolution-betti,b3:cor:betti-gap-termination,b3:prop:pd-minimal-length}给出的对应集中列于\cref{b3:tab:minimal-resolution-dictionary}。

\begin{center}
\begin{minipage}{\textwidth}
\makeatletter
\def\@captype{table}
\makeatother
\centering
\caption[极小自由消解的同调解释]{极小自由消解的同调解释。设 \((R,\mathfrak m)\) 为交换 Noether 局部环，\(k=R/\mathfrak m\)，\(M\ne0\) 为有限生成的 \(R\) 模，\(F_\bullet\rightarrow M\rightarrow0\) 为其极小自由消解。下表中 \(i\ge0\)，关系模指极小满射 \(F_0\rightarrow M\) 的核。}
\label{b3:tab:minimal-resolution-dictionary}
\begin{tabularx}{\textwidth}{@{}>{\raggedright\arraybackslash}p{0.36\textwidth}>{\raggedright\arraybackslash}X@{}}
\toprule
极小消解数据 & 同调解释或数值意义\\
\midrule
\(\operatorname{rank}_R F_i\) & \(\beta_i^R(M)\)\\
\(F_i/\mathfrak mF_i\) & \(\operatorname{Tor}_i^R(M,k)\)\\
\(\operatorname{Hom}_R(F_i,k)\) & \(\operatorname{Ext}_R^i(M,k)\)\\
\(F_0\) 的秩 & \(M\) 的最少生成元数\\
\(F_1\) 的秩 & 关系模的最少生成元数\\
某一阶 \(\beta_i^R(M)=0\) & 等价于 \(F_i=0\)、\(\operatorname{Tor}_i^R(M,k)=0\) 或 \(\operatorname{Ext}_R^i(M,k)=0\)，并迫使 \(F_j=0\) 对所有 \(j\ge i\) 成立\\
最后非零项为 \(F_r\) & \(\operatorname{pd}_RM=r\)\\
每个 \(F_i\) 都非零 & \(\operatorname{pd}_RM=\infty\)，没有最后非零项\\
\bottomrule
\end{tabularx}
\end{minipage}
\end{center}

\begin{proposition}\label{b3:prop:pd-test-module-counterexamples}
设 \(k\) 为域，\(R=k[x,y]_{(x,y)}\)、\(M=R/(x)\)。则
\[
\operatorname{Tor}_1^R(M,R/(y))=0,\qquad
\operatorname{Tor}_1^R(M,k)\cong k,\qquad \operatorname{pd}_RM=1.
\]
另令 \(S=k[\epsilon]/(\epsilon^2)\)，把 \(k=S/(\epsilon)\) 视为 \(S\) 模。则
\[
\begin{gathered}
\operatorname{Ext}_S^i(k,S)=0\quad(i>0),\qquad
\operatorname{Ext}_S^i(k,k)\cong k\quad(i\ge0),\\
\operatorname{pd}_Sk=\infty.
\end{gathered}
\]
因此投射维数的局部判据中，不能将剩余域任意换成别的测试模。
\end{proposition}
\begin{proof}
\(x\) 在 \(R\) 上为非零因子，故
\[
0\longrightarrow R\xrightarrow{x}R\longrightarrow M\longrightarrow0
\]
是自由消解。\(x\in(x,y)R\) 使它极小，最后非零自由项在次数一，所以 \(\operatorname{pd}_RM=1\)。张量 \(R/(y)\cong k[x]_{(x)}\) 后，乘法 \(x\) 仍单射，一次 Tor 为零；张量剩余域 \(k\) 后，微分却为零，一次 Tor 等于 \(k\)。非零的测试模 \(R/(y)\) 没有检测出这一项。

在 \(S\) 中，使用\secref{b3:sec:1701}给出的周期消解：每一项为 \(S\)，全部正次微分为乘 \(\epsilon\)。它们的核与像都为 \((\epsilon)\)，所以消解正合；各项秩为一，微分又属于极大理想，故它极小。由\cref{b3:prop:pd-minimal-length}，\(\operatorname{pd}_Sk=\infty\)。映入 \(S\) 后，上链复形的微分仍为乘 \(\epsilon\)，其正次的核与像相等，故正次 Ext 为零。映入 \(k\) 后，微分全部为零，每个次数的 Ext 都为 \(k\)。因此即使第二变量取环本身，也可能看不见无限消解中的全部正次项。
\end{proof}

\begin{corollary}\label{b3:cor:completion-betti-pd}
设 \((R,\mathfrak m)\) 为交换 Noether 局部环，\(k=R/\mathfrak m\) 为剩余域，\(M\) 为有限生成的 \(R\) 模。令 \(\widehat R\) 为 \(R\) 的 \(\mathfrak m\) 进完备化，\(\widehat M=\widehat R\otimes_RM\)。则对所有整数 \(i\ge0\)，有
\[
\beta_i^{\widehat R}(\widehat M)=\beta_i^R(M),
\qquad \operatorname{pd}_{\widehat R}\widehat M=\operatorname{pd}_RM.
\]
投射维数的等式包括 \(\infty\) 的情形；\(M=0\) 时两边均按约定为 \(-\infty\)，全部 Betti 数为零。
\end{corollary}
\begin{proof}
\(M=0\) 时直接成立，以下设 \(M\ne0\)。由\cref{b3:thm:completion-tensor}，所定义的 \(\widehat M\) 也是 \(M\) 的 \(\mathfrak m\) 进完备化。\cref{b3:thm:completion-flat}给出 \(R\to\widehat R\) 忠实平坦，所以 \(\widehat M\ne0\)；它又是有限生成的 \(\widehat R\) 模。由\cref{b3:prop:completed-local-ring}，\(\widehat R\) 是 Noether 局部环，其极大理想为 \(\mathfrak m\widehat R\)，剩余域自然等于 \(k\)。

取 \(M\) 的极小自由消解 \(F_\bullet\)。平坦性使 \(F_\bullet\otimes_R\widehat R\to\widehat M\to0\) 仍正合，各项是同秩的有限自由模。原微分的矩阵系数在 \(\mathfrak m\) 中，变基后仍在 \(\mathfrak m\widehat R\) 中，所以这也是极小自由消解。再张量共同的剩余域，得到
\[
(F_i\otimes_R\widehat R)\otimes_{\widehat R}k
\cong F_i\otimes_Rk.
\]
由\cref{b3:prop:minimal-resolution-betti}，两边的 \(k\) 维数就是各自的 Betti 数，故逐次相等。这两个极小消解的非零项出现在完全相同的次数；再用\cref{b3:prop:pd-minimal-length}，得到有限或无限的投射维数相等。
\end{proof}
